%
% Copyright (C) 1995, 2000, 2013 American Mathematical Society.
% Copyright (C) 2016-2025 LaTeX Project and American Mathematical Society.
%
% This work may be distributed and/or modified under the
% conditions of the LaTeX Project Public License, either version 1.3c
% of this license or (at your option) any later version.
% The latest version of this license is in
%   https://www.latex-project.org/lppl.txt
% and version 1.3c or later is part of all distributions of LaTeX
% version 2005/12/01 or later.
%
% This work has the LPPL maintenance status `maintained'.
%
% The Current Maintainer of this work is the LaTeX Project.
%
\documentclass[leqno,titlepage,openany]{amsldoc}[1999/12/13]

\def\MaintainedByLaTeXTeam#1{%
\begin{center}%
\fbox{\fbox{\begin{tabular}{@{}l@{}}%
This file is maintained by the \LaTeX{} Project team.\\%
Bug reports can be opened (category \texttt{#1}) at\\%
\url{https://latex-project.org/bugs/}.\end{tabular}}}\end{center}}

\ifx\UndEfiNed\url
  \ClassError{amsldoc}{%
    This version of amsldoc.tex must be processed\MessageBreak
    with a newer version of amsldoc.cls (2.02 or later).}{}
\fi

\renewcommand\bibname{References}

%    Definitions that should be moved to amsldoc.cls
%    when it is next updated
\makeatletter
\let\cleardouble@page\cleardoublepage
\AtBeginDocument{%
  \ifx\cleardouble@page\cleardoublepage
    \def\cleardoublepage{\clearpage{\pagestyle{empty}\cleardouble@page}}
  \fi
}

% #1 is empty or one of class, environment, option, package
% could use different font in each case
\makeatletter
\def\@category@index#1#2{%
    \autoindex{%
        #2\actualchar\string\index#1font{#2}%
        \ifx\@nil#1\@nil\else\space#1\fi
    }%
}

\let\indexpackagefont\textsf
\let\indexclassfont\textsf
\let\indexoptionfont\texttt
\let\indexenvironmentfont\texttt
\let\indexfont\texttt
\makeatother
%    For a package that won't be indexed, i.e. amsmath
\newcommand{\nipkg}{\textsf}
\def\tikz/{Ti\textit{k}Z}
%    End definitions for amsldoc.cls

\title{User's Guide for the \nipkg{amsmath} Package (Version~2.1)}
\author{American Mathematical Society, \LaTeX\ Project}
\date{1999-12-13\\(revised 2002-02-25, 2016-11-14, 2018-04-05, 2019-10-14, 2020-02-18)}
\makeatletter
\def\@thanks{\bigskip\MaintainedByLaTeXTeam{amslatex}}
\makeatother

% Very minor changes are not listed in the \date above:
%
% 2021-08-22: reported typo fixed mathools -> mathtools
% 2025-05-23: drop spurious option list


%    Use the amsmath package and amscd package in order to print
%    examples.
\usepackage{amsmath}
\usepackage{amscd}

\usepackage{imakeidx}
%    Need alternate coding for vertical bar to get it into index.
%    Based on \qcbang and \cnbang in amsldoc.cls, refined by David Carlisle.
%    Give it a sort field that will be at the end of the symbols;
%    must avoid " and | which have special meanings.
\DeclareRobustCommand{\qcvert}{\qc\|\index{{z1}@{\ntt\char'174}}}
\DeclareRobustCommand{\cnvert}{%
  \ncn{\|}\index{{z0}@{\ntt\bslchar\char'174}}}
%    Need alternate index sort field for bang to allow & to sort first
\DeclareRobustCommand{\altcnbang}{%
  \ncn{\!}\index{"!b@{\ntt\bslchar\qcbang}}}
\DeclareRobustCommand{\cnamp}{%
  \index{"!a@\texttt{\&}}}

\makeindex[intoc] % generate index data
\providecommand{\see}[2]{\textit{see} #1}

%    The amsldoc class includes a number of features useful for
%    documentation about TeX, including:
%
%    ---Commands \tex/, \amstex/, \latex/, ... for uniform treatment
%    of the various logos and easy handling of following spaces.
%
%    ---Commands for printing various common elements: \cn for command
%    names, \fn for file names (including font-file names), \env for
%    environments, \pkg and \cls for packages and classes, etc.

%    Many of the command names used here are rather long and will
%    contribute to poor linebreaking if we follow the \latex/ practice
%    of not hyphenating anything set in tt font; instead we selectively
%    allow some hyphenation.
\allowtthyphens % defined in amsldoc.cls

\hyphenation{ac-cent-ed-sym-bol add-to-counter add-to-length align-at
  aligned-at allow-dis-play-breaks ams-art ams-cd ams-la-tex amsl-doc
  ams-symb ams-tex ams-text ams-xtra bmatrix bold-sym-bol cen-ter-tags
  eqn-ar-ray idots-int int-lim-its latex med-space neg-med-space
  neg-thick-space neg-thin-space no-int-lim-its no-name-lim-its
  over-left-arrow over-left-right-arrow over-right-arrow pmatrix
  qed-sym-bol set-length side-set small-er tbinom the-equa-tion
  thick-space thin-space un-der-left-arrow un-der-left-right-arrow
  un-der-right-arrow use-pack-age var-inj-lim var-proj-lim vmatrix
  xalign-at xx-align-at}

%    Prepare for illustrating the \vec example
\newcommand{\vect}[1]{\mathbf{#1}}

\newcommand{\booktitle}[1]{\textit{#1}}
\newcommand{\journalname}[1]{\textit{#1}}
\newcommand{\seriesname}[1]{\textit{#1}}

%    Command to insert and index a particular phrase. Doesn't work for
%    certain kinds of special characters in the argument.
\newcommand{\ii}[1]{#1\index{#1}}

\newcommand{\vstrut}[1]{\vrule width0pt height#1\relax}

%    An environment for presenting comprehensive address information:
\newenvironment{infoaddress}{%
  \par\topsep\medskipamount
  \trivlist\centering
  \item[]%
  \begin{minipage}{.7\columnwidth}%
  \raggedright
}{%
  \end{minipage}%
  \endtrivlist
}

\newenvironment{eqxample}{%
  \par\addvspace\medskipamount
  \noindent\begin{minipage}{.5\columnwidth}%
  \def\producing{\end{minipage}\begin{minipage}{.5\columnwidth}%
    \hbox\bgroup\kern-.2pt\vrule width.2pt%
      \vbox\bgroup\parindent0pt\relax
%    The 3pt is to cancel the -\lineskip from \displ@y
    \abovedisplayskip3pt \abovedisplayshortskip\abovedisplayskip
    \belowdisplayskip0pt \belowdisplayshortskip\belowdisplayskip
    \noindent}
}{%
  \par
%    Ensure that a lonely \[\] structure doesn't take up width less than
%    \hsize.
  \hrule height0pt width\hsize
  \egroup\vrule width.2pt\kern-.2pt\egroup
  \end{minipage}%
  \par\addvspace\medskipamount
}

%    Errors in output shouldn't be \texttt, and they may also be longer
%    than one line, but they should be identified by the "errorbullet".
%    (They were formerly labeled as aubsections.)
\newenvironment{erroro}[1]{%
  \par\addvspace\medskipamount
  \noindent\hangindent\parindent
  \hbox to\parindent{\errorbullet\hfil}\ignorespaces
  #1\par\vspace{\smallskipamount}
}{%
  \par\addvspace\medskipamount
}

%    The chapters are so short, perhaps we shouldn't call them by the
%    name `Chapter'. We make \chaptername read an argument in order to
%    remove a following \space or "{} " (both possibilities are present
%    in book.cls).

\renewcommand{\chaptername}[1]{}
\newcommand{\chapnum}[1]{\mdash #1\mdash }
\makeatletter
\def\@makechapterhead#1{%
  \vspace{1.5\baselineskip}%
  {\parindent \z@ \raggedright \reset@font
    \ifnum \c@secnumdepth >\m@ne
      \large\bfseries \chapnum{\thechapter}%
      \par\nobreak
      \vskip.5\baselineskip\relax
    \fi
    #1\par\nobreak
    \vskip\baselineskip
  }}
\makeatother

%    A command for ragged-right parbox in a tabular.
\newcommand{\rp}{\let\PBS\\\raggedright\let\\\PBS}

%    Non-indexed file name
\newcommand{\nfn}[1]{\texttt{#1}}

%    For the examples in the math spacing table.
%%\newcommand{\lspx}{\mbox{\rule{5pt}{.6pt}\rule{.6pt}{6pt}}}
%%\newcommand{\rspx}{\mbox{\rule[-1pt]{.6pt}{7pt}%
%%  \rule[-1pt]{5pt}{.6pt}}}
\newcommand{\lspx}{\mathord{\Rightarrow\mkern-1mu}}
\newcommand{\rspx}{\mathord{\mkern-1mu\Leftarrow}}
\newcommand{\spx}[1]{$\lspx #1\rspx$}

%    For a list of characters representing document input.
\newcommand{\clist}[1]{%
  \mbox{\ntt\spaceskip.2em plus.1em \xspaceskip\spaceskip#1}}

%    Fix weird \latex/ definition of rightmark.
\makeatletter
\def\rightmark{\expandafter\@rightmark\botmark{}{}}
%    Also turn off section marks.
\let\sectionmark\@gobble
\renewcommand{\chaptermark}[1]{%
  \uppercase{\markboth{\rhcn#1}{\rhcn#1}}}
\newcommand{\rhcn}{\thechapter. }
\makeatother

%    Include down to \section but not \subsection, in toc:
\setcounter{tocdepth}{1}

%    Subheadings for the bibliography
\newif\iffirstbibsubhead \firstbibsubheadtrue
\newcommand{\bibsubhead}[1]{%
  \iffirstbibsubhead \firstbibsubheadfalse
  \else \addvspace{\medskipamount}
  \fi
  \item[]\hspace*{-\leftmargin}\textbf{#1}\par
  \vspace{\smallskipamount}
}

\DeclareMathOperator{\ix}{ix}
\DeclareMathOperator{\nul}{nul}
\DeclareMathOperator{\End}{End}
\DeclareMathOperator{\xxx}{xxx}

\usepackage{hyperref}
\hypersetup{pdfborder=0 0 0}
\makeatletter
\let\oldcs\cs
\def\cs#1{\texorpdfstring{\oldcs{#1}}{\@backslashchar\@backslashchar#1}}
\let\cn\cs
\makeatother
\begin{document}
\chapter{Integrals and sums}

\section{Multiline subscripts and superscripts}

The \cn{substack} command can be used to produce a multiline subscript
or superscript:\index{subscripts and superscripts!multiline}\relax
\index{superscripts|see{subscripts and superscripts}} for example
\begin{ctab}{ll}
\begin{minipage}[t]{.6\columnwidth}
\begin{verbatim}
\sum_{\substack{
         0\le i\le m\\
         0<j<n}}
  P(i,j)
\end{verbatim}
\end{minipage}
&
$\displaystyle
\sum_{\substack{0\le i\le m\\ 0<j<n}} P(i,j)$
\end{ctab}
A slightly more generalized form is the \env{subarray} environment which
allows you to specify that each line should be left-aligned instead of
centered, as here:
\begin{ctab}{ll}
\begin{minipage}[t]{.6\columnwidth}
\begin{verbatim}
\sum_{\begin{subarray}{l}
        i\in\Lambda\\ 0<j<n
      \end{subarray}}
 P(i,j)
\end{verbatim}
\end{minipage}
&
$\displaystyle
  \sum_{\begin{subarray}{l}
        i\in \Lambda\\ 0<j<n
      \end{subarray}}
 P(i,j)$
\end{ctab}

%% ------------------------------------------------------------------ %%

\section{The \cn{sideset} command}\label{sideset}

There's also a command called \cn{sideset}, for a rather special
purpose: putting symbols at the subscript and
superscript\index{subscripts and superscripts!on sums} corners of a
large operator symbol such as $\sum$ or $\prod$. \emph{Note: this
command is not designed to be applied to anything other than sum-class
symbols.} The prime
example is the case when you want to put a prime on a sum symbol. If
there are no limits above or below the sum, you could just use
\cn{nolimits}: here's
%%%%%%%%%%%%%%%%%%%%%%%%%%%%%%%%%%%%%%%%%%%%%%%%%%%%%%%%%%%%%%%%%%%%%%%%
|\sum\nolimits' E_n| in display mode:
\begin{equation}
\sum\nolimits' E_n
\end{equation}
If, however, you want not only the prime but also something below or
above the sum symbol, it's not so easy\mdash indeed, without
\cn{sideset}, it would be downright difficult. With \cn{sideset}, you
can write
\begin{ctab}{ll}
\begin{minipage}[t]{.6\columnwidth}
\begin{verbatim}
\sideset{}{'}
  \sum_{n<k,\;\text{$n$ odd}} nE_n
\end{verbatim}
\end{minipage}
&$\displaystyle
\sideset{}{'}\sum_{n<k,\;\text{$n$ odd}} nE_n
$
\end{ctab}
The extra pair of empty braces is explained by the fact that
\cn{sideset} has the capability of putting an extra symbol or symbols at
each corner of a large operator; to put an asterisk at each corner of a
product symbol, you would type
\begin{ctab}{ll}
\begin{minipage}[t]{.6\columnwidth}
\begin{verbatim}
\sideset{_*^*}{_*^*}\prod
\end{verbatim}
\end{minipage}
&$\displaystyle
\sideset{_*^*}{_*^*}\prod
$
\end{ctab}

%% ------------------------------------------------------------------ %%

\section{Placement of subscripts and limits}\label{subplace}

The default positioning for subscripts depends on the
base symbol involved. The default for sum-class symbols is
`displaylimits' positioning: When a sum-class symbol appears
in a displayed formula, subscript and superscript are placed in `limits'
position above and below, but in an inline formula, they are placed to
the side, to avoid unsightly and wasteful spreading of the
surrounding text lines.
The default for integral-class symbols is to have sub- and
superscripts always to the side, even in displayed formulas.
(See the discussion of the \opt{intlimits} and related
options\index{options!behavior of particular options} in
Section~\ref{options}.)

Operator names such as $\sin$ or $\lim$ may have either `displaylimits'
or `limits' positioning depending on how they were defined. The standard
operator names are defined according to normal mathematical usage.

The commands \cn{limits} and \cn{nolimits} can be used to override the
normal behavior of a base symbol:
\begin{equation*}
\sum\nolimits_X,\qquad \iint\limits_{A},
\qquad\varliminf\nolimits_{n\to \infty}
\end{equation*}
To define a command whose subscripts follow the
same `displaylimits' behavior as \cn{sum}, put
\cn{displaylimits} at the tail end of the definition. When multiple
instances of \cn{limits}, \cn{nolimits}, or \cn{displaylimits} occur
consecutively, the last one takes precedence.

%% ------------------------------------------------------------------ %%

\section{Multiple integral signs}

\cn{iint}, \cn{iiint}, and \cn{iiiint} give multiple integral
signs\index{integrals!multiple} with the spacing between them nicely
adjusted, in both text and display style. \cn{idotsint} is an extension
of the same idea that gives two integral signs with dots between them.
\begin{gather}
\iint\limits_A f(x,y)\,dx\,dy\qquad\iiint\limits_A
f(x,y,z)\,dx\,dy\,dz\\
\iiiint\limits_A
f(w,x,y,z)\,dw\,dx\,dy\,dz\qquad\idotsint\limits_A f(x_1,\dots,x_k)
\end{gather}

%%%%%%%%%%%%%%%%%%%%%%%%%%%%%%%%%%%%%%%%%%%%%%%%%%%%%%%%%%%%%%%%%%%%%%%%

\chapter{Commutative diagrams}\label{s:commdiag}
\index{commutative diagrams}

Some commutative diagram commands like the ones in \amstex/ are
available as a separate package, \pkg{amscd}. For complex commutative
diagrams authors will need to turn to more comprehensive packages like
\tikz/\index{TikZ@\tikz/ package} (in particular, \pkg{tikz-cd})
or \xypic/\index{XY-pic@\xypic/ package},
but for simple diagrams without diagonal
arrows\index{arrows!in commutative diagrams} the \pkg{amscd} commands
may be more convenient. Here is one example.
\begin{equation*}
\begin{CD}
S^{\mathcal{W}_\Lambda}\otimes T   @>j>>   T\\
@VVV                                    @VV{\End P}V\\
(S\otimes T)/I                  @=      (Z\otimes T)/J
\end{CD}
\end{equation*}
\begin{verbatim}
\begin{CD}
S^{\mathcal{W}_\Lambda}\otimes T   @>j>>   T\\
@VVV                                    @VV{\End P}V\\
(S\otimes T)/I                  @=      (Z\otimes T)/J
\end{CD}
\end{verbatim}
In the \env{CD} environment the commands |@>>>|,
|@<<<|, |@VVV|, and |@AAA| give respectively right, left, down, and up
arrows. For the horizontal arrows, material between the first and second
|>| or |<| symbols will be typeset as a superscript, and material
between the second and third will be typeset as a subscript. Similarly,
material between the first and second or second and third |A|s or |V|s
of vertical arrows will be typeset as left or right \qq{sidescripts}.
The commands |@=| and \verb'@|' give horizontal and vertical double lines.
A \qq{null arrow} command |@.| can be used instead of a visible arrow
to fill out an array where needed.

%%%%%%%%%%%%%%%%%%%%%%%%%%%%%%%%%%%%%%%%%%%%%%%%%%%%%%%%%%%%%%%%%%%%%%%%

\chapter{Using math fonts}

\section{Introduction}

For more comprehensive information on font use in \latex/, see the
\latex/ font guide (\fn{fntguide.tex}) or \booktitle{The \latex/
Companion} \cite{tlc2}. The basic set of math font commands\index{math
fonts}\relax \index{math symbols|see{math fonts}} in \latex/ includes
\cn{mathbf}, \cn{mathrm}, \cn{mathcal}, \cn{mathsf}, \cn{mathtt},
\cn{mathit}. Additional math alphabet commands such as
\cn{mathbb} for blackboard bold, \cn{mathfrak} for Fraktur, and
\cn{mathscr} for Euler script are available through the packages
\pkg{amsfonts} and \pkg{euscript} (distributed separately).

%% ------------------------------------------------------------------ %%
\enlargethispage{1\baselineskip}

\section{Recommended use of math font commands}

If you find yourself employing math font commands frequently in your
document, you might wish that they had shorter names, such as \ncn{mb}
instead of \cn{mathbf}. Of course, there is nothing to keep you from
providing such abbreviations for yourself by suitable \cn{newcommand}
statements. But for \latex/ to provide shorter names would actually be a
disservice to authors, as that would obscure a much better alternative:
defining custom command names derived from the names of the underlying
mathematical objects, rather than from the names of the fonts used to
distinguish the objects. For example, if you are using bold to indicate
vectors, then you will be better served in the long run if you define a
`vector' command instead of a `math-bold' command:
\begin{verbatim}
  \newcommand{\vect}[1]{\mathbf{#1}}
\end{verbatim}
you can write |\vect{a} + \vect{b}| to produce $\vect{a} +
\vect{b}$.
If you decide several months down the road that you want to use the bold
font for some other purpose, and mark vectors by a small over-arrow
instead, then you can put the change into effect merely by changing the
definition of \ncn{vect}; otherwise you would have to replace all
occurrences of \cn{mathbf} throughout your document, perhaps even
needing to inspect each one to see whether it is indeed
an instance of a vector.

It can also be useful to assign distinct
command names for different letters of a particular font:
\begin{verbatim}
\DeclareSymbolFont{AMSb}{U}{msb}{m}{n}% or use amsfonts package
\DeclareMathSymbol{\C}{\mathalpha}{AMSb}{"43}
\DeclareMathSymbol{\R}{\mathalpha}{AMSb}{"52}
\end{verbatim}
These statements would define the commands \cn{C} and \cn{R} to produce
blackboard-bold letters from the `AMSb' math symbols font. If you refer
often to the\break complex numbers or real numbers in your document, you
might find this method more convenient than (let's say) defining a
\ncn{field} command and writing\break |\field{C}|, |\field{R}|. But for
maximum flexibility and control, define such a \ncn{field} command and
then define \ncn{C} and \ncn{R} in terms of that
command:\indexcs{mathbb}
\begin{verbatim}
\usepackage{amsfonts}% to get the \mathbb alphabet
\newcommand{\field}[1]{\mathbb{#1}}
\newcommand{\C}{\field{C}}
\newcommand{\R}{\field{R}}
\end{verbatim}

%% ------------------------------------------------------------------ %%
\enlargethispage{3\baselineskip}

\section{Bold math symbols}

The \cn{mathbf} command is commonly used to obtain bold Latin letters in
math, but for most other kinds of math symbols it has no effect, or its
effects depend unreliably on the set of math fonts that are in use. For
example, writing
\begin{verbatim}
\Delta \mathbf{\Delta}\mathbf{+}\delta \mathbf{\delta}
\end{verbatim}
produces $\Delta \mathbf{\Delta}\mathbf{+}\delta
\mathbf{\delta}$; the \cn{mathbf} has no effect on the plus
sign or the small delta.
% should be paragraph break; avoid overfull box on next line
The \nipkg{amsmath} package therefore provides two additional commands,
\cn{boldsymbol} and \cn{pmb}, that can be applied to other kinds of math
symbols. \cn{boldsymbol} can be used for a math symbol that remains
unaffected by \cn{mathbf} if (and only if) your current math font set
includes a bold version of that symbol. \cn{pmb} can be used as a last
resort for any math symbols that do not have a true bold version
provided by your set of math fonts; \qq{pmb} stands for \qq{poor man's
bold} and the command works by typesetting multiple copies of the symbol
with slight offsets. The quality of the output is inferior, especially
for symbols that contain any hairline strokes. When the standard default set of
\latex/ math fonts are in use (Computer Modern), the only symbols that
are likely to require \cn{pmb} are large operator symbols like \cn{sum},
extended delimiter\index{delimiters} symbols, or the extra math symbols
provided by the \pkg{amssymb} package \cite{amsfonts}.

The following formula shows some of the results that are possible:
\begin{verbatim}
A_\infty + \pi A_0
\sim \mathbf{A}_{\boldsymbol{\infty}} \boldsymbol{+}
  \boldsymbol{\pi} \mathbf{A}_{\boldsymbol{0}}
\sim\pmb{A}_{\pmb{\infty}} \pmb{+}\pmb{\pi} \pmb{A}_{\pmb{0}}
\end{verbatim}
\begin{equation*}
A_\infty + \pi A_0
\sim \mathbf{A}_{\boldsymbol{\infty}} \boldsymbol{+}
  \boldsymbol{\pi} \mathbf{A}_{\boldsymbol{0}}
\sim\pmb{A}_{\pmb{\infty}} \pmb{+}\pmb{\pi} \pmb{A}_{\pmb{0}}
\end{equation*}
If you want to use only the \cn{boldsymbol} command without loading the
whole \nipkg{amsmath} package, the \pkg{bm} package is recommended (this
is a standard \latex/ package, not an AMS package; you probably have it
already if you have a 1997 or newer version of \latex/).

%% ------------------------------------------------------------------ %%

\section{Italic Greek letters}

For italic versions of the capital Greek letters, the following commands
are provided:
\begin{ctab}{rlrl}
\cn{varGamma}& $\varGamma$& \cn{varSigma}& $\varSigma$\\
\cn{varDelta}& $\varDelta$& \cn{varUpsilon}& $\varUpsilon$\\
\cn{varTheta}& $\varTheta$& \cn{varPhi}& $\varPhi$\\
\cn{varLambda}& $\varLambda$& \cn{varPsi}& $\varPsi$\\
\cn{varXi}& $\varXi$& \cn{varOmega}& $\varOmega$\\
\cn{varPi}& $\varPi$
\end{ctab}

%%%%%%%%%%%%%%%%%%%%%%%%%%%%%%%%%%%%%%%%%%%%%%%%%%%%%%%%%%%%%%%%%%%%%%%%

\chapter{Error messages and output problems}

\section{General remarks}

This is a supplement to Chapter~8 of the \latex/ manual \cite{lamport}
(first edition: Chapter~6). Appendix~B of the
\textit{Companion}~\cite{tlc2} lists all \latex/ errors including those
for \nipkg{amsmath} with likely causes for each.  For the reader's
convenience, the set of error messages discussed here overlaps
somewhat with the set in those references,
but please be aware that we don't provide exhaustive coverage here.
The error messages are arranged in alphabetical order, disregarding
unimportant text such as |! LaTeX Error:| at the beginning, and
nonalphabetical characters such as \qc{\\}. Where examples are given, we
show also the help messages that appear on screen when you respond to an
error message prompt by entering |h|.

Remember that the important error message is the first for a given line.
When running in nonstop mode, errors accumulate, and the first error
may create a further error condition from which it is impossible to
recover.  In such a situation processing will stop after 100~errors,
so the reporting of genuine errors may not be complete, and it may be
impossible to determine which reported errors are genuine and which
are not.

A final section discusses some output errors, i.e., instances where
the printed document has something wrong but there was no \latex/
error during typesetting.

%% ------------------------------------------------------------------ %%
%% **********
\enlargethispage{-2\baselineskip}

\end{document}
